% Einheit 1 – Lineares und exponentielles Wachstum unterscheiden
\setcounter{aufgabe}{9}
\hypertarget{e1}{}\einheitenkopf*{Einheit 1 von 5 · Lineares und exponentielles Wachstum unterscheiden}
\zweigzeile{Hier lernst du, an Tabellen, Texten und Graphen zu erkennen, ob ein Wert linear oder exponentiell wächst \,·\, neu in diesem Jahr \,·\, P10 oft \,·\, baut auf: Zunahme um denselben Betrag, Erhöhen um Prozent mit einer Dezimalzahl, Punkte im Koordinatensystem}

\begin{aufgabe}{Ich erkenne, ob sich ein Wert jedes Mal um denselben Betrag oder um denselben Prozentsatz ändert. Kreuze an. Du musst nichts rechnen.}
\begin{teile}
\teil Lea bekommt jeden Monat $3\,€$ mehr Taschengeld. \hfill \kreuz{gleicher Betrag}\kreuz{gleicher Prozentsatz}
\teil Ein Guthaben wächst jedes Jahr um $2\,\%$. \hfill \kreuz{gleicher Betrag}\kreuz{gleicher Prozentsatz}
\teil Eine Sonnenblume wird jede Woche $9\,\text{cm}$ größer. \hfill \kreuz{gleicher Betrag}\kreuz{gleicher Prozentsatz}
\teil Die Zahl der Follower steigt jeden Monat um $12\,\%$. \hfill \kreuz{gleicher Betrag}\kreuz{gleicher Prozentsatz}
\teil Der Wert eines Autos sinkt jedes Jahr um $15\,\%$. \hfill \kreuz{gleicher Betrag}\kreuz{gleicher Prozentsatz}
\end{teile}
\end{aufgabe}

\begin{aufgabe}{Ich erkenne, ob in einer Reihe immer dasselbe dazukommt oder immer mit derselben Zahl malgenommen wird. Kreuze an. Du musst nichts aufschreiben.}
\begin{teile}
\teil $2;\ 4;\ 8;\ 16$ \hfill \kreuz{gleiche Differenz}\kreuz{gleicher Quotient}
\teil $7;\ 10;\ 13;\ 16$ \hfill \kreuz{gleiche Differenz}\kreuz{gleicher Quotient}
\teil $81;\ 27;\ 9;\ 3$ \hfill \kreuz{gleiche Differenz}\kreuz{gleicher Quotient}
\teil $50;\ 42;\ 34;\ 26$ \hfill \kreuz{gleiche Differenz}\kreuz{gleicher Quotient}
\teil $4;\ 40;\ 400;\ 4\,000$ \hfill \kreuz{gleiche Differenz}\kreuz{gleicher Quotient}
\end{teile}
\end{aufgabe}

\begin{aufgabe}{Ich erkenne, wo ein Graph die senkrechte Achse trifft. Markiere den Punkt, an dem der Graph die senkrechte Achse trifft.}

\begin{ksys}[xmin=0,xmax=5,ymin=0,ymax=5,ablesen]
\funktion[3.5]{0.8*\x}{\mathrm{I}}
\end{ksys}\ksysabstand\begin{ksys}[xmin=0,xmax=5,ymin=0,ymax=5,ablesen]
\funktion[3.5]{1.2*1.3^\x}{\mathrm{II}}
\end{ksys}\ksysabstand\begin{ksys}[xmin=0,xmax=5,ymin=0,ymax=5,ablesen]
\funktion[3.5]{0.5*\x+2}{\mathrm{III}}
\end{ksys}

\begin{teile}
\teil bei Graph I \quad \teil bei Graph II \quad \teil bei Graph III
\teil Welcher Graph beginnt im Ursprung? \hfill \kreuz{I}\kreuz{II}\kreuz{III}
\end{teile}
\end{aufgabe}

\begin{aufgabe}{Ich kann an einer Tabelle entscheiden, ob ein Wert linear oder exponentiell wächst. Bilde die Differenzen aufeinanderfolgender Werte und entscheide, ob sie gleich sind.}
\beispiel{35;\ 43;\ 51;\ 59 &\qquad \text{Differenzen: } 8;\ 8;\ 8 &&\text{gleich, also lineares Wachstum}}
\begin{geruest}
\gz{a}{$12;\ 19;\ 26;\ 33$}{Differenzen: \leerfeld \quad gleich? \janein}
\gz{b}{$3;\ 13;\ 23;\ 33$}{Differenzen: \leerfeld \quad gleich? \janein}
\gz{c}{$10;\ 13;\ 17;\ 22$}{Differenzen: \leerfeld \quad gleich? \janein}
\gz{d}{$40;\ 55;\ 70;\ 85$}{Differenzen: \leerfeld \quad gleich? \janein}
\end{geruest}
Bilde die Quotienten aufeinanderfolgender Werte und entscheide, ob sie gleich sind.
\begin{geruest}
\gz{e}{$50;\ 55;\ 60{,}5;\ 66{,}55$}{Quotienten: \leerfeld \quad gleich? \janein}
\end{geruest}
Entscheide, ob der Wert linear oder exponentiell wächst, und begründe in einem Satz.
\begin{teile}
\teil Wert nach 0, 1, 2, 3 Jahren: $800;\ 840;\ 882;\ 926{,}1$ \hfill \kreuz{linear}\kreuz{exponentiell}
\teil Die Zahl der Nutzer einer App wächst jeden Monat um $15\,\%$. \hfill \kreuz{linear}\kreuz{exponentiell}
\teil Wert nach 0, 1, 2, 3 Jahren: $2\,000;\ 1\,800;\ 1\,620;\ 1\,458$. Entscheide auch, ob der Wert zunimmt oder abnimmt. \\ \kreuz{linear}\kreuz{exponentiell}\qquad\kreuz{nimmt zu}\kreuz{nimmt ab}
\teil Berechne für f) die Zuwächse von Jahr zu Jahr. Erkläre, warum sie größer werden, obwohl der Prozentsatz gleich bleibt.
\teil Die Fläche einer Wasserpflanze auf einem Teich verdoppelt sich jede Woche. Erkläre, warum das nicht immer so weitergehen kann.
\teil Ein Café notiert seinen Jahresumsatz: 2022: $240\,000\,€$, 2023: $252\,000\,€$, 2024: $264\,600\,€$, 2025: $277\,830\,€$. Entscheide, ob der Umsatz linear oder exponentiell wächst. Begründe deine Entscheidung. (P10 2018)
\end{teile}
\end{aufgabe}

\begin{aufgabe}{Ich kann einen Graphen zu einem Wachstum auswählen. Entscheide für jeden Graphen, ob er eine Gerade oder eine gekrümmte Kurve ist.}

\begin{ksys}[xmin=0,xmax=4,ymin=0,ymax=4,ablesen]
\funktion[2.5]{0.9*\x}{\mathrm{I}}
\end{ksys}\hspace{0.5cm}\begin{ksys}[xmin=0,xmax=4,ymin=0,ymax=4,ablesen]
\funktion[2.5]{1.2*1.35^\x}{\mathrm{II}}
\end{ksys}\hspace{0.5cm}\begin{ksys}[xmin=0,xmax=4,ymin=0,ymax=4,ablesen]
\funktion[2.5]{0.6*\x+1.5}{\mathrm{III}}
\end{ksys}\hspace{0.5cm}\begin{ksys}[xmin=0,xmax=4,ymin=0,ymax=4,ablesen]
\funktion[2.5]{0.25*\x^2}{\mathrm{IV}}
\end{ksys}

\begin{teilezwei}
\tz Graph I: \kreuz{Gerade}\kreuz{Kurve} & \tz Graph II: \kreuz{Gerade}\kreuz{Kurve} \\
\tz Graph III: \kreuz{Gerade}\kreuz{Kurve} & \tz Graph IV: \kreuz{Gerade}\kreuz{Kurve} \\
\end{teilezwei}
\begin{teile}
\teil Ein Bestand beginnt mit einem Anfangswert größer als null. Welche Graphen scheiden deshalb aus? \\ \kreuz{I}\kreuz{II}\kreuz{III}\kreuz{IV}
\teil Der Bestand wächst jedes Jahr um denselben Prozentsatz. Kreuze den passenden Graphen an. \\ \kreuz{I}\kreuz{II}\kreuz{III}\kreuz{IV}
\end{teile}
\end{aufgabe}

\begin{aufgabe}{Ich kann einen Graphen zu einem Wachstum auswählen – weiter. Ein E-Bike kostet neu $3\,200\,€$ und verliert jedes Jahr $15\,\%$ seines Werts. Waagerecht ist die Zeit in Jahren abgetragen, senkrecht der Wert in €.}

\begin{ksys}[xmin=0,xmax=6,ymin=0,ymax=3600,ystep=600,ablesen]
\funktion[3]{3200-300*\x}{\mathrm{I}}
\end{ksys}\hspace{0.6cm}\begin{ksys}[xmin=0,xmax=6,ymin=0,ymax=3600,ystep=600,ablesen]
\funktion[3]{3200*0.85^\x}{\mathrm{II}}
\end{ksys}\hspace{0.6cm}\begin{ksys}[xmin=0,xmax=6,ymin=0,ymax=3600,ystep=600,ablesen]
\funktion[3]{2000*0.85^\x}{\mathrm{III}}
\end{ksys}

\begin{teile}
\teil Begründe in einem Satz, warum die Gerade I nicht passt.
\teil Begründe, warum Graph III nicht passt.
\teil Begründe, warum Graph II passt: Woran erkennst du an seinem Verlauf die Abnahme um denselben Prozentsatz?
\end{teile}
Ein Streamingdienst hat $40\,000$ Abonnenten. Ihre Zahl steigt jeden Monat um $9\,\%$. Waagerecht ist die Zeit in Monaten abgetragen, senkrecht die Zahl der Abonnenten in Tausend.

\begin{ksys}[xmin=0,xmax=6,ymin=0,ymax=80,ystep=10,ablesen]
\funktion[3]{40+3.6*\x}{\mathrm{I}}
\end{ksys}\hspace{0.6cm}\begin{ksys}[xmin=0,xmax=6,ymin=0,ymax=80,ystep=10,ablesen]
\funktion[4]{60*(1.09^\x-1)}{\mathrm{II}}
\end{ksys}\hspace{0.6cm}\begin{ksys}[xmin=0,xmax=6,ymin=0,ymax=80,ystep=10,ablesen]
\funktion[3]{40*1.09^\x}{\mathrm{III}}
\end{ksys}

\begin{teile}
\teil Entscheide, welcher Graph die Zahl der Abonnenten beschreibt. Begründe für die beiden anderen Graphen, warum sie nicht passen. (P10 2026)
\end{teile}
\end{aufgabe}

\begin{aufgabe}{Ich kann Wertepaare eines Wachstums als Punkte in ein Koordinatensystem eintragen. Die Tabelle zeigt, wie viele Bakterien (in Tausend) in einer Probe gezählt wurden.}
\begin{minipage}[t]{0.52\linewidth}\vspace{0pt}
\begin{ksys}[xmin=0,xmax=4,ymin=0,ymax=90,xstep=0.5,ystep=10,xlabel=$t$ in h,ylabel=Tausend]
\end{ksys}
\end{minipage}\hfill
\begin{minipage}[t]{0.46\linewidth}\vspace{0pt}
\sachtabelle{lcccc}{Zeit in h & 0 & 1 & 2 & 3}{Bakterien in Tausend & 10 & 20 & 40 & 80}
Trage den Punkt für den angegebenen Zeitpunkt ein.
\begin{teilezwei}
\tz $t = 0$ & \tz $t = 1$ \\
\tz $t = 2$ & \tz $t = 3$ \\
\end{teilezwei}
\begin{teile}
\teil Nach $2{,}5$ Stunden wurden $57$ Tausend Bakterien gezählt. Trage auch diesen Punkt so genau wie möglich ein.
\teil Verbinde die Punkte zu einer Kurve.
\end{teile}
\end{minipage}
\end{aufgabe}

\begin{aufgabe}{Ich kann Wertepaare eines Wachstums als Punkte in ein Koordinatensystem eintragen – weiter. Wähle selbst eine Einteilung für beide Achsen und beschrifte sie.}
\begin{teile}
\teil Ein Sparplan: Nach 0, 2, 4, 6 und 8 Jahren beträgt das Guthaben $500\,€$, $605\,€$, $732\,€$, $886\,€$ und $1\,072\,€$. Trage die Punkte ein (linkes Gitter).
\teil Die Helligkeit unter Wasser nimmt mit jedem Meter Tiefe ab. In 0, 1, 2 und 3 Metern Tiefe misst man $800$, $520$, $338$ und $220$ Lux. Wähle eine geeignete Einteilung der Achsen, trage die vier Punkte ein und verbinde sie zu einer Kurve (rechtes Gitter). (P10 2017)
\end{teile}

\leeresgitter{9}{11}\hspace{0.8cm}\leeresgitter{9}{11}

\end{aufgabe}

\begin{aufgabe}{Ich finde den Fehler bei der Entscheidung zwischen linear und exponentiell. Der Preis eines Sammlerstücks beträgt nach 0, 1, 2, 3 Jahren $500\,€$, $560\,€$, $627{,}20\,€$, $702{,}46\,€$; er steigt jedes Jahr um $12\,\%$. Jonas schreibt:}
\par\hspace{1.6em}\parbox{0.9\linewidth}{„Der Preis steigt jedes Jahr um $12\,\%$, also immer um gleich viel. Das ist lineares Wachstum.“}\par
\begin{teile}
\teil Finde den Fehler, benenne ihn und korrigiere Jonas' Entscheidung.
\teil Entscheide selbst und begründe: Die Einwohnerzahl eines Dorfes sinkt jedes Jahr um $2\,\%$: $35\,000$; $34\,300$; $33\,614$; $32\,942$. \hfill \kreuz{linear}\kreuz{exponentiell}
\end{teile}
\end{aufgabe}

\begin{aufgabe}{Ich kann begründen, warum gleicher Prozentsatz nicht gleicher Betrag heißt. Entscheide ohne Rechnung und begründe.}
\begin{teile}
\teil Ein Guthaben wächst jedes Jahr um $2\,\%$. Tim sagt: „Dann bekomme ich jedes Jahr gleich viel Zinsen.“ Hat Tim recht?
\teil Ein Wert nimmt jedes Jahr um $10\,\%$ ab. Werden die Verluste in Euro von Jahr zu Jahr größer, kleiner, oder bleiben sie gleich?
\teil Ein Bestand wächst exponentiell und hat den Anfangswert $500$. Kann sein Graph durch den Ursprung gehen?
\end{teile}
\end{aufgabe}

\begin{aufgabe}{Ich kann mit linearem und exponentiellem Wachstum eine Entscheidung treffen. Ein Radverleih hat $1\,200$ Räder. Plan A: Er kauft jedes Jahr $150$ Räder dazu. Plan B: Er vergrößert den Bestand jedes Jahr um $10\,\%$.}
\sachtabelle{lccccccccc}{Jahr & 0 & 1 & 2 & 3 & 4 & 5 & 6 & 7 & 8}{Plan A & 1\,200 & 1\,350 & 1\,500 & 1\,650 & 1\,800 & 1\,950 & 2\,100 & 2\,250 & 2\,400\\ Plan B & 1\,200 & 1\,320 & 1\,452 & 1\,597 & 1\,757 & 1\,933 & 2\,126 & 2\,338 & 2\,572}
\begin{teile}
\teil Welcher Plan beschreibt lineares, welcher exponentielles Wachstum? Begründe mit der Tabelle.
\teil Ab welchem Jahr hat der Verleih mit Plan B mehr Räder als mit Plan A?
\teil Der Verleih hat nur Platz für $2\,300$ Räder. Mit welchem Plan ist der Platz zuerst zu klein, und in welchem Jahr?
\end{teile}
\end{aufgabe}
